\subsection{A Stylized View of Finite-Budget Detectability}
\label{sec:minervarl-support-bridging}

Let $a^\star$ denote the ground-truth structured answer for prompt $x$. We define full verifier success directly through the task reward,
\[
    S(x,y;a^\star)
    :=
    \mathbb{1}\!\left\{
    R_{\textsc{minerva}}(x,y,a^\star)=1
    \right\},
\]
and let
\[
    p_\theta(a^\star\mid x)
    :=
    \Pr_{y\sim\pi_\theta(\cdot\mid x)}
    \left[S(x,y;a^\star)=1\right].
\]
This definition uses the verifier rather than equality of extracted answers, thereby covering normalization, aliases, set-valued outputs, and graded task rewards. With $k$ independent rollouts, the probability that the rollout group contains no fully verified completion is
\begin{equation}
    \Pr[\text{no full success in $k$ rollouts}]
    =
    \left(1-p_\theta(a^\star\mid x)\right)^k .
\end{equation}
For a failure tolerance $\zeta\in(0,1)$, define the detectability threshold
\begin{equation}
    \varepsilon_{k,\zeta}
    :=
    1-\zeta^{1/k}
    \approx
    \frac{-\log \zeta}{k}.
\end{equation}
If $p_\theta(a^\star\mid x)<\varepsilon_{k,\zeta}$, then a fully verified response is not reliably observed under the available rollout budget.

This full-success event should not be conflated with an all-zero reward group. Writing $m=\max_j R_{\textsc{minerva}}(x,y_j,a^\star)$, MinervaRL's ACR gate uses $m<1$, whereas the zero-solve statistic in \Cref{sec:reward-sparsity-dynamics} uses $m=0$. When $0<m<1$, no rollout is fully verified, but unequal partial-credit rewards can still give GRPO nonzero relative advantages. The calculation above concerns only whether full verifier success is detectable; it does not characterize all learning signal carried by partial rewards.

For a gated prompt, the answer-conditioned prompt $x^{\mathrm{acr}}=x\oplus b(a^\star,d)$ makes it more likely that the EMA teacher will generate a fully verified trace. After verifier and quality filtering, an accepted trace is distilled onto the original prompt $x$, without a label hint. This step is intended to raise the original-prompt probability of full success so that later on-policy rollouts can observe it under the same budget $k$.

The corresponding appendix result is deliberately conditional and stylized. It assumes both that ACR generation yields an accepted trace with probability at least $\alpha$ and that each effective original-prompt distillation update raises the log probability of full verifier success by at least $\Delta$. Given those assumptions, it bounds the number of accepted updates needed to cross the finite-budget detectability threshold. It does not establish that the assumed increase occurs, nor does it model joint GRPO--SFT optimization, gradient interference across prompts, or partial-credit reward dynamics. We therefore use the result only as finite-sampling intuition for the method, not as an optimization or convergence theory of MinervaRL. Appendix~\ref{app:minervarl-support-bridging} gives the formal statement and its scope.
